\chapter*{Введение}
\addcontentsline{toc}{chapter}{Введение}
\markboth{Введение}{}

Язык Java существует больше тридцати лет, и всё это время его хоронили. Сначала он
был «слишком медленным», потом «слишком многословным», потом «устаревшим на фоне
новых языков». Тем не менее сегодня на Java написана значительная часть банковских
систем, торговых площадок, государственных информационных систем и корпоративных
приложений~--- то есть ровно того программного обеспечения, которое должно работать
годами и не ломаться при обновлении. Причина проста: Java оказалась удачным
компромиссом между скоростью разработки, производительностью и предсказуемостью
сопровождения.

Это пособие~--- о том, как писать на Java не учебные примеры, а работающие
приложения. Разница между этими двумя занятиями больше, чем кажется. Написать класс
с двумя полями умеет второкурсник; собрать веб-сервис, который переживёт
одновременное обращение сотни пользователей, корректно откатит транзакцию при сбое
и не развалится при первом же изменении требований,~--- это другой навык. Именно ему
и посвящена книга.

\section*{Как устроено пособие}
\addcontentsline{toc}{section}{Как устроено пособие}

Материал разбит на двенадцать глав, выстроенных так, чтобы каждая опиралась на
предыдущие.

\emph{Главы 1--5} закладывают фундамент: устройство платформы Java и виртуальной
машины, синтаксис языка, объектно-ориентированное программирование, обработка
исключений и отладка, коллекции, ввод-вывод и многопоточность. Это тот минимум,
без которого дальнейшее чтение бессмысленно.

\emph{Главы 6--7} переводят разговор на уровень реальных проектов: системы сборки,
работа с базами данных через JDBC и Hibernate, функциональный стиль обработки
данных и, наконец, Spring~--- фреймворк, на котором сегодня написано большинство
серверных приложений на Java.

\emph{Главы 8--12} посвящены тому, что окружает код: веб-архитектуре и протоколу
HTTP, графическим интерфейсам на JavaFX, документированию и тестированию,
паттернам проектирования и системам контроля версий вместе с процессами
разработки.

Каждая глава построена одинаково. Сначала идёт теория с разбором примеров, затем
\textbf{практикум}~--- задания, которые нужно выполнить за компьютером, потом
\textbf{контрольные вопросы} для самопроверки и \textbf{тестовые задания} с
вариантами ответов. Правильные ответы на тесты вынесены в приложение~\ref{app:keys},
чтобы соблазн подсмотреть возникал хотя бы не сразу.

\section*{Что нужно, чтобы читать эту книгу}
\addcontentsline{toc}{section}{Что нужно, чтобы читать эту книгу}

Предполагается, что читатель уже писал программы на каком-нибудь языке и понимает,
что такое переменная, цикл и функция. Знание именно Java не требуется~--- первые
главы начинают с нуля.

Из инструментов понадобится компьютер с установленным JDK версии~21 или новее и
любая среда разработки. Как всё это установить и проверить, описано в
приложении~\ref{app:setup}. Специально оговорим: пособие рассчитано на работу
\emph{без доступа в интернет}. Весь код примеров приведён полностью, все команды
выписаны, все зависимости перечислены. Если для выполнения задания нужен файл
конфигурации~--- он есть в тексте, а не по ссылке.

\section*{Как работать с примерами}
\addcontentsline{toc}{section}{Как работать с примерами}

Код в книге стоит набирать руками, а не переписывать глазами. Это скучный совет,
но он работает: опечатка, из-за которой программа не компилируется, учит читать
сообщения об ошибках лучше любого объяснения. По той же причине в тексте
намеренно разобраны типичные ошибки~--- те, на которые студенты тратят часы,
не понимая, что происходит.

Листинги пронумерованы по главам, и на каждый есть ссылка в тексте. Длинные
фрагменты сокращены: если в примере важен один метод, остальное заменено
многоточием в комментарии. Такие сокращения всегда отмечены явно~--- если
многоточия нет, код можно набирать как есть и он заработает.

\section*{Результаты обучения}
\addcontentsline{toc}{section}{Результаты обучения}

Проработав пособие, читатель сможет:

\begin{itemize}
  \item объяснить, как устроена платформа Java и что происходит с программой
        от исходного текста до выполнения на виртуальной машине;
  \item уверенно пользоваться конструкциями языка, коллекциями и средствами
        обработки ошибок;
  \item проектировать классы, применяя принципы ООП и типовые паттерны, и
        распознавать антипаттерны в чужом коде;
  \item подключать базу данных через JDBC и Hibernate, понимая, что такое
        транзакция и уровень изоляции;
  \item собрать веб-приложение на Spring Boot с REST-интерфейсом,
        разграничением доступа и слоем представления;
  \item написать графический интерфейс на JavaFX;
  \item покрыть код тестами на JUnit~5 и Mockito и оформить документацию Javadoc;
  \item вести разработку в Git, разрешать конфликты слияния и настраивать
        автоматическую сборку.
\end{itemize}

\vspace{4mm}

Java~--- язык, который вознаграждает за аккуратность. Надеюсь, эта книга поможет
её приобрести.
